\documentclass{article}

% packages
\usepackage{amsmath, amsthm, thmtools, amsfonts, amssymb, luacode, catchfile, tikzducks, hyperref, ifthen}
\ifcsname c@kobocompile\endcsname
	\usepackage[a5paper, total={1072pt, 1448pt}, margin=10pt, includeheadfoot]{geometry} % set page margins
\else
	\usepackage[a4paper, margin=50pt, includeheadfoot]{geometry}
\fi
\usepackage[shortlabels]{enumitem}
\usepackage[skip=3pt, indent=0pt]{parskip}

% language
\usepackage[bidi=basic, layout=tabular, provide=*]{babel}
\ifcsname c@english\endcsname
	\babelprovide[main, import]{english}
\else
	\babelprovide[main, import]{hebrew}
	\babelprovide{rl}
\fi
%\babelfont{rm}{Libertinus Serif}
\babelfont{rm}[Renderer=Harfbuzz]{Libertinus Serif}
\babelfont{sf}{Libertinus Sans}
\babelfont{tt}{Libertinus Mono}

% style
\AddToHook{cmd/section/before}{\clearpage}	% Add line break before section
\linespread{1.3}
\setcounter{secnumdepth}{0}		% Remove default number tags from sections, this won't do well with theorems
\AtBeginDocument{\setlength{\belowdisplayskip}{3pt}}
\AtBeginDocument{\setlength{\abovedisplayskip}{3pt}}
\graphicspath{ {../images/} }

% operators
\DeclareMathOperator\cis{cis}
\DeclareMathOperator\Sp{Sp}
\DeclareMathOperator\tr{tr}
\DeclareMathOperator\im{Im}
\DeclareMathOperator\re{Re}
\DeclareMathOperator\diag{diag}
\DeclareMathOperator*\lowlim{\underline{lim}}
\DeclareMathOperator*\uplim{\overline{lim}}
\DeclareMathOperator\rng{rng}
\DeclareMathOperator\Sym{Sym}
\DeclareMathOperator\Arg{Arg}
\DeclareMathOperator\Log{Log}
\DeclareMathOperator\dom{dom}
\DeclareMathOperator\supp{Supp}
\DeclareMathOperator\var{Var}
\DeclareMathOperator\cov{Cov}

% commands
%\renewcommand\qedsymbol{\textbf{מש''ל}}
%\renewcommand\qedsymbol{\fbox{\emoji{lizard}}}
\newcommand{\Aa}[0]{\mathcal{A}}
\newcommand{\Bb}[0]{\mathcal{B}}
\newcommand{\CC}[0]{\mathbb{C}}
\newcommand{\Cc}[0]{\mathcal{C}}
\newcommand{\EE}[0]{\mathbb{E}}
\newcommand{\FF}[0]{\mathbb{F}}
\newcommand{\Ff}[0]{\mathcal{F}}
\newcommand{\Ii}[0]{\mathcal{I}}
\newcommand{\Gg}[0]{\mathcal{G}}
\newcommand{\Ll}[0]{\mathcal{L}}
\newcommand{\Mm}[0]{\mathcal{M}}
\newcommand{\NN}[0]{\mathbb{N}}
\newcommand{\Nn}[0]{\mathcal{N}}
\newcommand{\PP}[0]{\mathbb{P}}
\newcommand{\Pp}[0]{\mathcal{P}}
\newcommand{\QQ}[0]{\mathbb{Q}}
\newcommand{\RR}[0]{\mathbb{R}}
\newcommand{\Rr}[0]{\mathcal{R}}
\newcommand{\Ss}[0]{\mathcal{S}}
\newcommand{\TT}[0]{\mathbb{T}}
\newcommand{\Uu}[0]{\mathcal{U}}
\newcommand{\Vv}[0]{\mathcal{V}}
\newcommand{\Ww}[0]{\mathcal{W}}
\newcommand{\ZZ}[0]{\mathbb{Z}}
\newcommand{\acts}[0]{\circlearrowright}
\newcommand{\explain}[2] {
	\begin{flalign*}
		 && \text{#2} && \text{#1}
	\end{flalign*}
}
\newcommand{\maketitleprint}[0]{ \begin{center}
	%\begin{tikzpicture}[scale=3]
	%	\duck[graduate=gray!20!black, tassel=red!70!black]
	%\end{tikzpicture}	
	\includegraphics[width=6cm]{cover}
\end{center}
}

% theorem commands
\newtheoremstyle{c_remark}
	{}	% Space above
	{}	% Space below
	{}% Body font
	{}	% Indent amount
	{\bfseries}	% Theorem head font
	{}	% Punctuation after theorem head
	{.5em}	% Space after theorem head
	{\thmname{#1}\thmnumber{ #2}\thmnote{ \normalfont{\text{(#3)}}}}	% head content
\newtheoremstyle{c_definition}
	{3pt}	% Space above
	{3pt}	% Space below
	{}% Body font
	{}	% Indent amount
	{\bfseries}	% Theorem head font
	{}	% Punctuation after theorem head
	{.5em}	% Space after theorem head
	{\thmname{#1}\thmnumber{ #2}\thmnote{ \normalfont{\text{(#3)}}}}	% head content
\newtheoremstyle{c_plain}
	{3pt}	% Space above
	{3pt}	% Space below
	{\itshape}% Body font
	{}	% Indent amount
	{\bfseries}	% Theorem head font
	{}	% Punctuation after theorem head
	{.5em}	% Space after theorem head
	{\thmname{#1}\thmnumber{ #2}\thmnote{ \text{(#3)}}}	% head content

\ifcsname c@english\endcsname
	\theoremstyle{plain}
	\newtheorem{theorem}{Theorem}[section]
	\newtheorem{lemma}[theorem]{Lemma}
	\newtheorem{proposition}[theorem]{Proposition}
	\newtheorem*{proposition*}{Proposition}
	%\newtheorem{corollary}[theorem]{אין חלופה עברית}

	\theoremstyle{definition}
	\newtheorem{definition}[theorem]{Definition}
	\newtheorem*{definition*}{Definition}
	\newtheorem{example}{Example}[section]
	\newtheorem{exercise}{Exercise}[section]

	\theoremstyle{remark}
	\newtheorem*{remark}{Remark}
	\newtheorem*{solution}{Solution}
	\newtheorem{conclusion}[theorem]{Conclusion}
	\newtheorem{notation}[theorem]{Notation}
\else
	\theoremstyle{c_plain}
	\newtheorem{theorem}{משפט}[section]
	\newtheorem{lemma}[theorem]{למה}
	\newtheorem{proposition}[theorem]{טענה}
	\newtheorem*{proposition*}{טענה}
	%\newtheorem{corollary}[theorem]{אין חלופה עברית}

	\theoremstyle{c_definition}
	\newtheorem{definition}[theorem]{הגדרה}
	\newtheorem*{definition*}{הגדרה}
	\newtheorem{example}{דוגמה}[section]
	\newtheorem{exercise}{תרגיל}[section]

	\theoremstyle{c_remark}
	\newtheorem*{remark}{הערה}
	\newtheorem*{solution}{פתרון}
	\newtheorem{conclusion}[theorem]{מסקנה}
	\newtheorem{notation}[theorem]{סימון}
\fi

% Questions related commands
\newcounter{question}
\setcounter{question}{1}
\newcounter{sub_question}
\setcounter{sub_question}{1}

\ifcsname c@english\endcsname
	\newcommand{\question}[1][0]{
		\ifthenelse{#1 = 0}{}{\setcounter{question}{#1}}
		\section{Question \arabic{question}}
		\addtocounter{question}{1}
		\setcounter{sub_question}{1}
	}

	\newcommand{\subquestion}[1][0]{
		\ifthenelse{#1 = 0}{}{\setcounter{sub_question}{#1}}
		\subsection{Part \alph{sub_question}}
		\addtocounter{sub_question}{1}
	}
\else
	\newcommand{\question}[1][0]{
		\ifthenelse{#1 = 0}{}{\setcounter{question}{#1}}
		\section{שאלה \arabic{question}}
		\addtocounter{question}{1}
		\setcounter{sub_question}{1}
	}

	\newcommand{\subquestion}[1][0]{
		\ifthenelse{#1 = 0}{}{\setcounter{sub_question}{#1}}
		\subsection{סעיף \localecounter{letters.gershayim}{sub_question}}
		\addtocounter{sub_question}{1}
	}
\fi

% import lua and start of document
\directlua{common = require ('../common')}

\GetEnv{AUTHOR}

% headers
\author{\AUTHOR}
\date\today

\title{פתרון מטלה 04 --- מבוא לטופולוגיה, 80516}
% chktex-file 9
% chktex-file 17

\begin{document}
\maketitle
\maketitleprint[purple]

\question{}
יהי $X$ מרחב טופולוגי ותהי $A \subseteq X$ קשירה.
נראה שכל $B$ המקיימת $A \subseteq B \subseteq \overline{A}$ היא קשירה.
\begin{proof}
	נניח בשלילה ש־$B$ לא קשירה, ויהיו $B \subseteq B_0 \uplus B_1$ קבוצות פתוחות זרות המעידות על כך.
	בלי הגבלת הכלליות נניח ש־$A \subseteq B_0$ (לא יתכן שגם $A \cap B_1 \ne \emptyset$).
	$B_1$ פתוחה ולכן $B_1^C$ סגורה ובפרט $\overline{A} \setminus B_1$ קבוצה סגורה המכילה את $A$, אבל $\overline{A}$ מינימלית בהכלתה את $A$,
	לכן נובע $\overline{A} = \overline{A} \setminus B_1$, כלומר $B_1 \cap \overline{A} = \emptyset$.
	נסיק אם כן ש־$B_1 = \emptyset$ בלבד, אבל $B_0, B_1$ לא ריקות מהגדרת הקשירות, ולכן קיבלנו סתירה.
\end{proof}

\question[3]
ניזכר ונגדיר שמרחב מטרי חסום אם הוא מוכל בכדור.

\subquestion{}
יהי $(X, \rho)$ מרחב מטרי קומפקטי,
נראה שהוא חסום.
\begin{proof}
	יהי $\epsilon > 0$ ונגדיר $U_x = B(x, \epsilon)$ לכל $x \in X$.
	אז $\bigcup_{x \in X} U_x = X$ הוא כיסוי פתוח, ולכן מקומפקטיות יש תת־כיסוי סופי $X = \bigcup_{i = 1}^n U_x$ עבור $x_1, \ldots, x_n \in X$ ו־$n \in \NN$.
	נגדיר $d = \sup_{i \le n} \rho(x_1, x_i)$ ונבחן את $B(x_1, r + \epsilon)$.
	כל נקודה $x \in X$ מקיימת $x \in B(x_i, \epsilon)$ עבור איזשהו $i$, וכן $\rho(x_1, x) \le \rho(x_1, x_i) + \rho(x_i, x) \le r + \epsilon$ ולכן $x \in B(x_1, r + \epsilon)$, כלומר כדור זה חוסם את $X$.
\end{proof}

\subquestion{}
נמצא דוגמה למרחב מטרי חסום שאינו קומפקטי.
\begin{solution}
	נתחיל בהבחנה שאנו יודעים, במרחבים מטריים קומפקטיות שקולה לקומפקטיות סדרתית, לכן מספיק שנמצא מרחב מטרי חסום ולא קומפקטי סדרתי.
	נגדיר $X = (0, 1)$ יחד עם המטריקה הסטנדרטית $\rho(x, y) = |x - y|$, ונבחין כי $X = B(\frac{1}{2}, \frac{1}{2})$.
	נגדיר ${\{ x_n \}}_{n = 1}^\infty \subseteq X$ על־ידי $x_n = \frac{1}{n}$,
	ב־$\RR$ נקבל ש־$\lim_{n \to \infty} x_n = 0$, לכן $\{ x_n \}$ לא מתכנסת ב־$X$, ואילו יש לה תת־סדרה מתכנסת, אז היא מתכנסת ל־$0$ ולכן לא מתכנסת ב־$X$ גם, כלומר הסדרה מעידה על אי־קומפקטיות סדרתית.
	נסיק אם כך ש־$(X, \rho)$ הוא מרחב מטרי חסום אך לא קומפקטי.
\end{solution}

\question{}
נראה שאם $X$ אינסופית אז טופולוגיית המשלים בן־המנייה,
כלומר $A \subseteq X$ פתוחה אם ורק אם $|X \setminus A| \le \omega$,
היא לא קומפקטית.
\begin{proof}
	נסמן $|X| = \kappa$ ונגדיר $f : \kappa \to X$ כדי להימנע משימוש בבחירה.
	וידוע ש־$\kappa \ge \omega$ מהגדרת $\omega$ כסודר אינסופי מינימלי.
	לכן נוכל להסיק גם ש־$\kappa = \kappa + \omega$, כלומר קיימת תת־קבוצה $Y \subseteq X$ כך ש־$|X \setminus Y| = \omega$ וכן $|Y| = \kappa$.
	נגדיר $f : \omega \to X \setminus Y$ העתקה חד־חד ערכית ועל,
	ונגדיר $U_n = Y \cup \{ f(n) \}$ לכל $n < \omega$.
	נבחין כי $|X \setminus U_n| = |(X \setminus Y) \setminus \{ f(n) \}| = |X \setminus Y| = \omega$ ולכן $U_n$ קבוצה פתוחה לכל $n < \omega$.
	נובע אם כך,
	\[
		\bigcup_{n < \omega} U_n
		= \bigcup_{n < \omega} Y \cup \{ f(n) \}
		= Y \cup \bigcup_{n < \omega} \{ f(n) \}
		= Y \cup f(\NN)
		= X
	\]
	כלומר זהו כיסוי פתוח של $X$.
	לכל תת־קבוצה סופית $I \subseteq \omega$ נבחר $m = S(\sup I)$ ונקבל $f(m) \notin \bigcup_{n \in I} U_n$, כלומר מצאנו איבר שלא נמצא בתת־כיסוי סופי המושרה מ־$I$ לכל $I$ שנבחר.
	נסיק ש־$X$ לא קומפקטית.
\end{proof}

\question{}
יהי $X$ מרחב טופולוגי קומפקטי ותהי ${\{ x_n \}}_{n = 0}^\infty \subseteq X$ סדרה,
נסמן $A_N = \{ x_n \mid N \le n \in \NN \}$.
נראה ש־$\bigcap_{N = 0}^\infty \overline{A_N} \ne \emptyset$.
\begin{proof}
	מהתנאי השקול לקומפקטיות, לכל אוסף קבוצות סגורות של $X$ אשר מקיים את תכונת החיתוך הסופי, יש חיתוך לא ריק.
	נרצה אם כך להראות שתכונת החיתוך הסופי חלה.
	נניח ש־$I \subseteq \NN$ כך ש־$|I| < \omega$.
	נגדיר $M = \sup I$, מסופיות נובע ש־$M$ מתקבל, וכן מהגדרת $A_N$ נקבל שלכל $N \in I$ גם $N \le M$ ולכן $x_M \in A_N$.
	נובע ש־$x_M \in \bigcap_{i \in I} \overline{A_i}$, כלומר החיתוך שלהם לא ריק, ובהתאם תכונת החיתוך הסופי אכן נשמרת.
	נסיק אם כך שגם $\bigcap_{N = 0}^\infty \overline{A_N} \ne \emptyset$.
\end{proof}

\question{}
תהי קבוצה $S$ לא ריקה, ו־$D \subseteq \Pp(S)$ משפחת קבוצות המקיימת את תכונת החיתוך הסופי.
נוכיח שקיים מסנן $\langle D \rangle$ כך ש־$D \subseteq \langle D \rangle$ וכך שאם $F$ מסנן המקיים $D \subseteq F$, אז גם $\langle D \rangle \subseteq F$.
\begin{proof}
	נתחיל בהוכחה עבור המקרה $D = \emptyset$.
	במקרה זה נגדיר $\langle D \rangle = \{ S \}$, ונבחין כי אוסף זה מקיים $S \in \langle D \rangle, \emptyset \notin \langle D \rangle$, וכן מקיים באופן ריק סגירות לחיתוך ולקבוצות מכילות.
	כל מסנן $D \subseteq F \subseteq \Pp(S)$ מקיים $S \in F$ ולכן $\langle D \rangle \subseteq F$ ולכן סיימנו.

	נניח ש־$D \ne \emptyset$.
	נגדיר את הקבוצה $\langle D \rangle$ כסגירות לחיתוך ולהוספת איברים מעל $D$.
	כלומר נגדיר $D \subseteq \langle D \rangle$ וכן לכל $A, B \in \langle D \rangle$ גם $A \cap B \in \langle D \rangle$ וכן לכל $\langle D \rangle \ni A \subseteq B \subseteq S$ גם $B \in \langle D \rangle$.
	נראה ש־$\langle D \rangle$ מסנן.
	ידוע ש־$D$ לא ריקה ולכן קיים $A \in D$, אבל $A \subseteq S$ ולכן $S \in \langle D \rangle$.
	אם $\emptyset \in D$ אז $D$ לא מקיימת את תכונת החיתוך הסופי, לכן נסיק ש־$\emptyset \notin D$, וכן לכל $A, B \in D$ ידוע ש־$A \cap B \ne \emptyset$, ואף לכל $A \in \langle D \rangle$ לא מתקיים $A \subsetneq \emptyset$,
	לכן נוכל להסיק ש־$\emptyset \notin \langle D \rangle$.
	הקבוצה $\langle D \rangle$ סגורה לחיתוך סופי מהגדרתה ככזו, וכן סגורה ללקיחת קבוצות מכילות גם מהגדרה, לכן נסיק ש־$\langle D \rangle$ אכן מסנן. \\
	נניח ש־$D \subseteq F$ מסנן כלשהו מעל $S$.
	לכל $A, B \in D$ גם $A, B \in F$ ולכן $A \cap B \in F$, וכן לכל $A \in D$ ו־$A \subseteq B$ גם $B \in F$, ולכן לכל $A \in \langle D \rangle$ נסיק ש־$A \in F$, כלומר $\langle D \rangle \subseteq F$.
\end{proof}

\question{}
תהי קבוצה $S$ ויהי $F \subseteq \Pp(S)$ מסנן על $S$.
יהיו $R$ קבוצה נוספת ו־$f : S \to R$ פונקציה כלשהי.
נגדיר $f_* F = \{ A \subseteq R \mid f^{-1}(A) \in F \}$, ונוכיח כי זהו מסנן מעל $R$.
\begin{proof}
	$f^{-1}(R) = S \in F$ מהגדרת $f$, לכן נובע ש־$R \in f_* F$.
	נבחין גם כי $\emptyset \notin F$ ולכן $f^{-1}(\emptyset) = \emptyset \notin f_* F$. \\
	נניח ש־$A, B \in f_* F$, אז $f^{-1}(A) \cap f^{-1}(B) \in F$ מהגדרת $F$.
	יהי $s \in f^{-1}(A) \cap f^{-1}(B)$, אז $f(s) \in A, B$ ונובע שבפרט גם $f(s) \in A \cap B$.
	נקבל אם כך ש־$s \in f^{-1}(f(s)) \subseteq f^{-1}(A \cap B)$, ולכן $f^{-1}(A) \cap f^{-1}(B) \subseteq f^{-1}(A \cap B)$ ולכן $f^{-1}(A \cap B) \in F$ ו־$A \cap B \in f_* F$. \\
	נניח ש־$A \in f_* F$, בהתאם $f^{-1}(A) \in F$, ונניח ש־$A \subseteq B$, אז גם $f^{-1}(A) \subseteq f^{-1}(B) \in F$ ולכן $B \in f_* F$. \\
	מצאנו אם כן ש־$f_* F$ מסנן מעל $R$.
\end{proof}

\end{document}
